% answers.tex —— 2013 年·福建卷·理科 · 参考答案与详细解析片段
% 由 tools/generate_projects.py 依来源总账自 2013/2013_数学真题_福建理_试题.md + 2013/2013_数学真题_福建理_解析正文.md + 2013/2013_数学真题_福建理_答案要点.md 生成，勿手改（改源 MD 后重跑）


\begin{answer}
\gksolhead{1}{}

\textbf{题目}：

已知复数 $z$ 的共轭复数 $\bar z=1+2\mathrm i$（$\mathrm i$ 为虚数单位），则 $z$ 在复平面内对应的点位于\mbox{（　　）}

（A）第一象限；（B）第二象限；（C）第三象限；（D）第四象限。

\Needspace{4\baselineskip}
\textbf{解析}

\Needspace{5\baselineskip}
由共轭复数的定义，若
\nopagebreak[4]
\[\bar z=1+2\mathrm i,\]
\Needspace{5\baselineskip}
则
\nopagebreak[4]
\[z=\overline{1+2\mathrm i}=1-2\mathrm i.\]
因此 $z$ 在复平面内对应的点为 $(1,-2)$，其横坐标为正、纵坐标为负，位于第四象限。

选 D。

\gkanswer{答案：D。$\bar z=1+2\mathrm i$，故 $z=1-2\mathrm i$，对应点 $(1,-2)$ 在第四象限。}

\end{answer}
\begin{answer}
\gksolhead{2}{}

\textbf{题目}：

已知集合 $A=\{1,a\}$，$B=\{1,2,3\}$，则“$a=3$”是“$A\subseteq B$”的\mbox{（　　）}

（A）充分而不必要条件；（B）必要而不充分条件；（C）充分必要条件；（D）既不充分也不必要条件。

\Needspace{4\baselineskip}
\textbf{解析}

\Needspace{5\baselineskip}
先看充分性。若 $a=3$，则
\nopagebreak[4]
\[A=\{1,3\}\subseteq\{1,2,3\}=B,\]
故“$a=3$”能推出“$A\subseteq B$”。

\Needspace{5\baselineskip}
再看必要性。若 $A\subseteq B$，由于 $A=\{1,a\}$，只须 $a\in B$，即
\nopagebreak[4]
\[a\in\{1,2,3\}.\]
取 $a=2$，此时 $A=\{1,2\}\subseteq B$，但 $a\neq3$。所以“$A\subseteq B$”不能推出“$a=3$”。

因此“$a=3$”是“$A\subseteq B$”的充分而不必要条件。

选 A。

\gkanswer{答案：A。$a=3$ 时 $A=\{1,3\}\subseteq B$；反之 $A\subseteq B$ 只要求 $a\in\{1,2,3\}$，取 $a=2$ 即不成立，故为充分而不必要条件。}

\end{answer}
\begin{answer}
\gksolhead{3}{}

\textbf{题目}：

双曲线 $\dfrac{x^2}{4}-y^2=1$ 的顶点到渐近线的距离等于\mbox{（　　）}

（A）$\dfrac25$；（B）$\dfrac45$；（C）$\dfrac{2\sqrt5}{5}$；（D）$\dfrac{4\sqrt5}{5}$。

\Needspace{4\baselineskip}
\textbf{解析}

\Needspace{5\baselineskip}
双曲线
\nopagebreak[4]
\[\frac{x^2}{4}-y^2=1\]
\Needspace{5\baselineskip}
中 $a=2$，$b=1$，顶点为 $(\pm2,0)$，渐近线为
\nopagebreak[4]
\[y=\pm\frac{x}{2},\]
\Needspace{5\baselineskip}
即
\nopagebreak[4]
\[x-2y=0,\qquad x+2y=0.\]
\Needspace{5\baselineskip}
取顶点 $(2,0)$，它到渐近线 $x-2y=0$ 的距离为
\nopagebreak[4]
\[d=\frac{|2-0|}{\sqrt{1^2+(-2)^2}}=\frac{2}{\sqrt5}=\frac{2\sqrt5}{5}.\]
由对称性，另一顶点到两条渐近线的距离也为同一值。

选 C。

\gkanswer{答案：C。双曲线顶点为 $(\pm2,0)$，渐近线为 $x\pm2y=0$，距离 $d=\dfrac{2}{\sqrt5}=\dfrac{2\sqrt5}{5}$。}

\end{answer}
\begin{answer}
\gksolhead{4}{}

\textbf{题目}：

某校从高一年级学生中随机抽取部分学生，将他们的模块测试成绩分成 $6$ 组：$[40,50)$，$[50,60)$，$[60,70)$，$[70,80)$，$[80,90)$，$[90,100)$ 加以统计，得到如图所示的频率分布直方图。已知高一年级共有学生 $600$ 名，据此估计，该模块测试成绩不少于 $60$ 分的学生人数为\mbox{（　　）}
\nopagebreak[4]
\begin{center}% ============================================================
% f2013_fujian_science_q04_histogram.tex
% 2013 福建卷（理工类）第 4 题 频率分布直方图（TikZ 矢量重绘）
% 组距 10，各组矩形高度按原页：0.005, 0.010, 0.015, 0.030, 0.025, 0.010
% 横轴 40--100 分；60 分以上四组高度之和为 0.080
% ============================================================
\begin{tikzpicture}[
  x=1cm, y=1cm, line width=0.6pt, font=\footnotesize, >=Stealth
]

  % ---------- 坐标轴 ----------
  \draw[->] (0,0) -- (8.9,0) node[below left, yshift=2pt] {分数};
  \draw[->] (0,0) -- (0,2.85) node[above] {};
  \node[align=center] at (0.58,3.10) {频率\\[-2pt]组距};
  \draw (0.18,3.08) -- (0.98,3.08);

  % 横轴断点符号
  \draw (0.32,0) -- (0.40,0.10) -- (0.48,0) -- (0.56,0.10) -- (0.64,0);

  % ---------- 纵向刻度与参考虚线 ----------
  \foreach \h/\lab in {0.005/0.005,0.010/0.010,0.015/0.015,0.025/0.025,0.030/0.030} {
    \pgfmathsetmacro{\yy}{\h*80}
    \draw[dashed, line width=0.35pt, gray!70] (0.08,\yy) -- (8.25,\yy);
    \node[left] at (-0.05,\yy) {\lab};
    \draw (-0.05,\yy) -- (0.05,\yy);
  }
  \node[left] at (-0.05,0) {$0$};

  % ---------- 各矩形 ----------
  % 40--50: 0.005
  \draw (3.20,0) rectangle (4.00,0.40);
  % 50--60: 0.010
  \draw (4.00,0) rectangle (4.80,0.80);
  % 60--70: 0.015
  \draw (4.80,0) rectangle (5.60,1.20);
  % 70--80: 0.030
  \draw (5.60,0) rectangle (6.40,2.40);
  % 80--90: 0.025
  \draw (6.40,0) rectangle (7.20,2.00);
  % 90--100: 0.010
  \draw (7.20,0) rectangle (8.00,0.80);

  % ---------- 横轴刻度 ----------
  \foreach \s/\xx in {40/3.20,50/4.00,60/4.80,70/5.60,80/6.40,90/7.20,100/8.00} {
    \draw (\xx,0) -- (\xx,-0.08);
    \node[below] at (\xx,-0.10) {\s};
  }

\end{tikzpicture}
\end{center}
\nopagebreak[4]
（A）$588$；（B）$480$；（C）$450$；（D）$120$。

\Needspace{4\baselineskip}
\textbf{解析}

\Needspace{5\baselineskip}
频率分布直方图中，每个矩形的面积等于该组的频率。组距为 $10$，故 $60$ 分以上四组的频率为
\nopagebreak[4]
\[\bigl(0.015+0.030+0.025+0.010\bigr)\times10=0.080\times10=0.80.\]
\Needspace{5\baselineskip}
所以成绩不少于 $60$ 分的学生人数约为
\nopagebreak[4]
\[600\times0.80=480.\]

选 B。

\begin{gktip}
说明：原图六组矩形高度依次为 $0.005,0.010,0.015,0.030,0.025,0.010$，六组之和为 $0.095$，与频率总和 $1$ 相差 $0.005$，疑为原图 $60$ 分以下某组印高有误。本题只读取 $60$ 分以上四组，读数与选项一致。
\end{gktip}

\gkanswer{答案：B。$60$ 分以上四组矩形高之和为 $0.015+0.030+0.025+0.010=0.080$，故人数为 $600\times0.080\times10=480$。}

\end{answer}
\begin{answer}
\gksolhead{5}{}

\textbf{题目}：

满足 $a,b\in\{-1,0,1,2\}$，且关于 $x$ 的方程 $ax^2+2x+b=0$ 有实数解的有序数对 $(a,b)$ 的个数为\mbox{（　　）}

（A）$14$；（B）$13$；（C）$12$；（D）$10$。

\Needspace{4\baselineskip}
\textbf{解析}

\Needspace{5\baselineskip}
方程
\nopagebreak[4]
\[ax^2+2x+b=0\]
中 $a,b\in\{-1,0,1,2\}$，按 $a$ 是否为 $0$ 分类。

\Needspace{5\baselineskip}
当 $a=0$ 时，方程化为
\nopagebreak[4]
\[2x+b=0,\]
\Needspace{5\baselineskip}
即 $x=-\dfrac b2$，对任意 $b\in\{-1,0,1,2\}$ 都有实数解，共
\nopagebreak[4]
\[4\text{ \text{对}}.\]

\Needspace{5\baselineskip}
当 $a\neq0$ 时，方程为一元二次方程，有实数解当且仅当判别式非负：
\nopagebreak[4]
\[\Delta=2^2-4ab=4(1-ab)\geqslant0\Longleftrightarrow ab\leqslant1.\]

逐个考察 $a\neq0$ 的三个取值：

\begin{itemize}
\item $a=-1$：$ab=-b\leqslant1$ 对 $b=-1,0,1,2$ 均成立，共 $4$ 对；
\item $a=1$：$ab=b\leqslant1$，得 $b=-1,0,1$，共 $3$ 对；
\item $a=2$：$ab=2b\leqslant1$，得 $b=-1,0$，共 $2$ 对。
\end{itemize}

\Needspace{5\baselineskip}
合计
\nopagebreak[4]
\[4+4+3+2=13\text{ \text{对}}.\]
注意 $a=0$ 时的 $4$ 对已单独计入，且不与上面重复。

选 B。

\gkanswer{答案：B。$a=0$ 时恒有实根，共 $4$ 对；$a\neq0$ 时由 $\Delta=4-4ab\geqslant0$ 得 $ab\leqslant1$，逐个计数为 $4+3+2=9$ 对，合计 $13$ 对。}

\end{answer}
\begin{answer}
\gksolhead{6}{}

\textbf{题目}：

阅读如图所示的程序框图，若输入的 $k=10$，则该算法的功能是\mbox{（　　）}
\nopagebreak[4]
\begin{center}% ============================================================
% f2013_fujian_science_q06_flowchart.tex
% 2013 福建卷（理工类）第 6 题 程序框图（TikZ 矢量重绘）
% S=0, i=1；输入 k；S=1+2S；i=i+1；i>k？否则回代，是则输出 S
% ============================================================
\begin{tikzpicture}[
  x=1cm, y=1cm, line width=0.6pt, font=\footnotesize, >=Stealth,
  every node/.style={inner sep=3pt},
  proc/.style={rectangle, draw, minimum width=2.6cm, minimum height=0.62cm},
  io/.style={trapezium, draw, trapezium left angle=70, trapezium right angle=110,
             minimum width=2.8cm, minimum height=0.66cm},
  dec/.style={diamond, draw, aspect=2.2, minimum width=2.8cm, minimum height=0.85cm},
  term/.style={rectangle, rounded corners=8pt, draw, minimum width=2.0cm, minimum height=0.58cm}
]

  \node[term] (start) at (0,0) {开始};
  \node[proc] (init)  at (0,-1.35) {$S=0,\ i=1$};
  \node[io]   (input) at (0,-2.60) {输入 $k$};
  \node[proc] (upd)   at (0,-3.85) {$S=1+2S$};
  \node[proc] (inc)   at (0,-5.10) {$i=i+1$};
  \node[dec]  (judge) at (0,-6.55) {$i>k$?};
  \node[io]   (out)   at (0,-8.10) {输出 $S$};
  \node[term] (end)   at (0,-9.35) {结束};

  \draw[->] (start) -- (init);
  \draw[->] (init) -- (input);
  \draw[->] (input) -- (upd);
  \draw[->] (upd) -- (inc);
  \draw[->] (inc) -- (judge);
  \draw[->] (judge) -- node[right] {是} (out);
  \draw[->] (out) -- (end);

  % 否：回到 S=1+2S 之前
  \draw[->] (judge.east) -- (2.60,-6.55) -- (2.60,-3.85) -- (upd.east);
  \node[above] at (1.30,-6.45) {否};

\end{tikzpicture}
\end{center}
\nopagebreak[4]
（A）计算数列 $\{2^{n-1}\}$ 的前 $10$ 项和；（B）计算数列 $\{2^{n-1}\}$ 的前 $9$ 项和；（C）计算数列 $\{2^n-1\}$ 的前 $10$ 项和；（D）计算数列 $\{2^n-1\}$ 的前 $9$ 项和。

\Needspace{4\baselineskip}
\textbf{解析}

\Needspace{5\baselineskip}
初始值 $S=0$，$i=1$。循环体先执行
\nopagebreak[4]
\[S\leftarrow1+2S,\qquad i\leftarrow i+1,\]
再判断 $i>k$。若未超过 $k$，返回继续执行循环体。

\Needspace{5\baselineskip}
设第 $n$ 次执行循环体后得到的值为 $S_n$，则
\nopagebreak[4]
\[S_n=1+2S_{n-1},\qquad S_0=0.\]
\Needspace{5\baselineskip}
于是
\nopagebreak[4]
\[S_1=1,\quad S_2=3,\quad S_3=7,\quad\cdots,\quad S_n=2^n-1.\]
\Needspace{5\baselineskip}
当 $k=10$ 时，循环体执行 $10$ 次后 $i=11>10$，输出
\nopagebreak[4]
\[S_{10}=2^{10}-1=1+2+2^2+\cdots+2^9.\]
这正是数列 $\{2^{n-1}\}$ 的前 $10$ 项和。

选 A。

\gkanswer{答案：A。循环执行 $10$ 次，$S=1+2+\cdots+2^9=2^{10}-1$，即数列 $\{2^{n-1}\}$ 的前 $10$ 项和。}

\end{answer}
\begin{answer}
\gksolhead{7}{}

\textbf{题目}：

在四边形 $ABCD$ 中，$\overrightarrow{AC}=(1,2)$，$\overrightarrow{BD}=(-4,2)$，则该四边形的面积为\mbox{（　　）}

（A）$\sqrt5$；（B）$2\sqrt5$；（C）$5$；（D）$10$。

\Needspace{4\baselineskip}
\textbf{解析}

\Needspace{5\baselineskip}
四边形的两条对角线为 $\overrightarrow{AC}$ 与 $\overrightarrow{BD}$。对任意四边形，面积等于两条对角线向量叉积绝对值的一半：
\nopagebreak[4]
\[S=\frac12\left|\overrightarrow{AC}\times\overrightarrow{BD}\right|.\]
\Needspace{5\baselineskip}
代入
\nopagebreak[4]
\[\overrightarrow{AC}=(1,2),\qquad\overrightarrow{BD}=(-4,2),\]
\Needspace{5\baselineskip}
得
\nopagebreak[4]
\[\overrightarrow{AC}\times\overrightarrow{BD}=1\times2-2\times(-4)=2+8=10.\]
\Needspace{5\baselineskip}
故
\nopagebreak[4]
\[S=\frac12\times10=5.\]

选 C。

\gkanswer{答案：C。四边形面积 $S=\dfrac12\left|\overrightarrow{AC}\times\overrightarrow{BD}\right|=\dfrac12|1\times2-2\times(-4)|=5$。}

\end{answer}
\begin{answer}
\gksolhead{8}{}

\textbf{题目}：

设函数 $f(x)$ 的定义域为 $\mathbf R$，$x_0$（$x_0\neq0$）是 $f(x)$ 的极大值点，以下结论一定正确的是\mbox{（　　）}

（A）$\forall x\in\mathbf R$，$f(x)\leqslant f(x_0)$；（B）$-x_0$ 是 $f(-x)$ 的极小值点；（C）$-x_0$ 是 $-f(x)$ 的极小值点；（D）$-x_0$ 是 $-f(-x)$ 的极小值点。

\Needspace{4\baselineskip}
\textbf{解析}

\Needspace{5\baselineskip}
已知 $x_0\neq0$ 是 $f(x)$ 的极大值点，即在 $x_0$ 的某邻域内恒有
\nopagebreak[4]
\[f(x)\leqslant f(x_0).\]

\Needspace{5\baselineskip}
考察 $g(x)=-f(-x)$。当 $x$ 在 $-x_0$ 附近时，$-x$ 在 $x_0$ 附近，故
\nopagebreak[4]
\[f(-x)\leqslant f(-x_0),\]
\Needspace{5\baselineskip}
两边乘以 $-1$ 得
\nopagebreak[4]
\[g(x)=-f(-x)\geqslant-f(-x_0)=g(-x_0).\]
所以 $-x_0$ 是 $g(x)=-f(-x)$ 的极小值点，D 正确。

其余选项：

\begin{itemize}
\item A 把“极大值”误当成“最大值”，极大值只是局部性质，不一定对全定义域成立；
\item B 中 $f(-x)$ 在 $-x_0$ 处应为极大值点，不是极小值点；
\item C 中 $-f(x)$ 在 $x_0$ 处为极小值点，与 $-x_0$ 无关。
\end{itemize}

选 D。

\gkanswer{答案：D。令 $g(x)=-f(-x)$。$f$ 在 $x_0$ 取极大值，故 $g$ 在 $-x_0$ 取极小值；其余三项均可举反例否定。}

\end{answer}
\begin{answer}
\gksolhead{9}{}

\textbf{题目}：

\Needspace{5\baselineskip}
已知等比数列 $\{a_n\}$ 的公比为 $q$，记
\nopagebreak[4]
\[b_n=a_{m(n-1)+1}+a_{m(n-1)+2}+\cdots+a_{m(n-1)+m},\qquad c_n=a_{m(n-1)+1}\cdot a_{m(n-1)+2}\cdot\cdots\cdot a_{m(n-1)+m},\]
其中 $m,n\in\mathbf N^*$，则以下结论一定正确的是\mbox{（　　）}

（A）数列 $\{b_n\}$ 为等差数列，公差为 $q^m$；（B）数列 $\{b_n\}$ 为等比数列，公比为 $q^{2m}$；（C）数列 $\{c_n\}$ 为等比数列，公比为 $q^{m^2}$；（D）数列 $\{c_n\}$ 为等比数列，公比为 $q^{m^m}$。

\Needspace{4\baselineskip}
\textbf{解析}

\Needspace{5\baselineskip}
设等比数列首项为 $a_1$，公比为 $q$。先看 $b_n$：
\nopagebreak[4]
\[b_n=a_{m(n-1)+1}+a_{m(n-1)+2}+\cdots+a_{m(n-1)+m}.\]
\Needspace{5\baselineskip}
当 $q\neq1$ 时，
\nopagebreak[4]
\[b_n=a_1q^{m(n-1)}\left(1+q+\cdots+q^{m-1}\right)
=a_1q^{m(n-1)}\cdot\frac{1-q^m}{1-q}.\]
\Needspace{5\baselineskip}
因此
\nopagebreak[4]
\[\frac{b_{n+1}}{b_n}=q^m,\]
$\{b_n\}$ 是公比为 $q^m$ 的等比数列，不是公差为 $q^m$ 的等差数列，也不是公比为 $q^{2m}$ 的等比数列。当 $q=1$ 时 $b_n=ma_1$ 为常数列，公比 $1=q^m$，同样符合。

\Needspace{5\baselineskip}
再看 $c_n$：
\nopagebreak[4]
\[c_n=a_{m(n-1)+1}a_{m(n-1)+2}\cdots a_{m(n-1)+m}.\]
\Needspace{5\baselineskip}
同样先设 $q\neq1$，则
\nopagebreak[4]
\[c_n=a_1^m q^{m\cdot m(n-1)}\cdot q^{0+1+\cdots+(m-1)}
=a_1^m q^{m^2(n-1)+\frac{m(m-1)}2}.\]
\Needspace{5\baselineskip}
于是
\nopagebreak[4]
\[\frac{c_{n+1}}{c_n}=q^{m^2}.\]
故 $\{c_n\}$ 是公比为 $q^{m^2}$ 的等比数列。当 $q=1$ 时 $c_n=a_1^m$ 为常数列，公比 $1=q^{m^2}$，仍符合。

选 C。

\gkanswer{答案：C。$c_{n+1}/c_n=q^{m^2}$，故 $\{c_n\}$ 是公比为 $q^{m^2}$ 的等比数列；$\{b_n\}$ 的公比应为 $q^m$，不是 $q^{2m}$。}

\end{answer}
\begin{answer}
\gksolhead{10}{}

\textbf{题目}：

设 $S$，$T$ 是 $\mathbf R$ 的两个非空子集，如果存在一个从 $S$ 到 $T$ 的函数 $y=f(x)$ 满足：

① $T=\{f(x)\mid x\in S\}$；

② 对任意 $x_1,x_2\in S$，当 $x_1<x_2$ 时，恒有 $f(x_1)<f(x_2)$；

那么称这两个集合“保序同构”。以下集合对不是“保序同构”的是\mbox{（　　）}

（A）$A=\mathbf N^*$，$B=\mathbf N$；（B）$A=\{x\mid-1\leqslant x\leqslant3\}$，$B=\{x\mid x=-8\text{ \text{或} }0<x\leqslant10\}$；（C）$A=\{x\mid0<x<1\}$，$B=\mathbf R$；（D）$A=\mathbf Z$，$B=\mathbf Q$。

\Needspace{4\baselineskip}
\textbf{解析}

“保序同构”要求存在从 $S$ 到 $T$ 的严格递增双射。

\begin{itemize}
\item A：$A=\mathbf N^*=\{1,2,3,\cdots\}$，$B=\mathbf N=\{0,1,2,\cdots\}$。令 $f(n)=n-1$，它是 $A$ 到 $B$ 的严格递增双射，故保序同构。
\item B：$A=[-1,3]$，$B=\{-8\}\cup(0,10]$。令
\end{itemize}
\nopagebreak[4]
\[f(-1)=-8,\qquad f(x)=\frac{10}{4}(x+1)\quad(-1<x\leqslant3).\]
  当 $x\to-1^+$ 时 $f(x)\to0^+$，当 $x=3$ 时 $f(3)=10$，故 $f$ 把 $(-1,3]$ 严格递增地映到 $(0,10]$，从而是 $A$ 到 $B$ 的严格递增双射。
\begin{itemize}
\item C：$A=(0,1)$，$B=\mathbf R$。令 $f(x)=\tan\left(\pi\left(x-\dfrac12\right)\right)$，它是 $(0,1)$ 到 $\mathbf R$ 的严格递增双射。
\item D：$A=\mathbf Z$，$B=\mathbf Q$。在 $\mathbf Z$ 中，每个整数 $n$ 都有“紧邻的后继” $n+1$，即不存在整数严格介于 $n$ 与 $n+1$ 之间；而 $\mathbf Q$ 中任意两个有理数之间都有有理数，没有这种紧邻后继。严格递增双射会保持“有紧邻后继”这一序性质，故 $\mathbf Z$ 与 $\mathbf Q$ 不可能保序同构。
\end{itemize}

选 D。

\gkanswer{答案：D。$\mathbf Z$ 中每个元素都有紧邻的后继，而 $\mathbf Q$ 中任意两数之间还有有理数，保序同构不可能把“有紧邻后继”这一序性质抹去；前三项均可构造严格递增的双射。}

\end{answer}
\begin{answer}
\gksolhead{11}{}

\textbf{题目}：

利用计算机产生 $0\sim1$ 之间的均匀随机数 $a$，则事件“$3a-1>0$”发生的概率为\gkblank{}。

\Needspace{4\baselineskip}
\textbf{解析}

\Needspace{5\baselineskip}
$a$ 在 $[0,1]$ 上均匀分布，事件
\nopagebreak[4]
\[3a-1>0\Longleftrightarrow a>\frac13.\]
\Needspace{5\baselineskip}
故所求概率为区间 $\left(\dfrac13,1\right]$ 的长度：
\nopagebreak[4]
\[P\left(a>\frac13\right)=1-\frac13=\frac23.\]

填 $\dfrac23$。

\gkanswer{答案：$\dfrac23$。$3a-1>0\Longleftrightarrow a>\dfrac13$，故概率为 $1-\dfrac13=\dfrac23$。}

\end{answer}
\begin{answer}
\gksolhead{12}{}

\textbf{题目}：

已知某一多面体内接于球构成一个简单组合体，如果该组合体的正视图、侧视图、俯视图均如图所示，且图中的四边形是边长为 $2$ 的正方形，则该球的表面积是\gkblank{}。
\nopagebreak[4]
\begin{center}% ============================================================
% f2013_fujian_science_q12_views.tex
% 2013 福建卷（理工类）第 12 题 三视图（TikZ 矢量重绘）
% 正视图、侧视图、俯视图均为：圆内含正方形（虚线），边长 2
% 正方体边长 2，外接球半径 sqrt(3)，表面积 12*pi
% ============================================================
\begin{tikzpicture}[
  x=1cm, y=1cm, line width=0.6pt, font=\footnotesize, >=Stealth
]

  \foreach \dx/\lab in {0/正视图, 4.05/侧视图, 8.10/俯视图} {
    \begin{scope}[xshift=\dx cm]
      % 外接圆
      \draw (0,0) circle (1.55);
      % 正方形（虚线）
      \draw[dashed] (-1,-1) rectangle (1,1);
      % 边长标注
      \draw[<->] (-1,-1.22) -- (1,-1.22);
      \node[below] at (0,-1.24) {$2$};
      % 视图名称
      \node at (0,-2.05) {\lab};
    \end{scope}
  }

\end{tikzpicture}
\end{center}
\nopagebreak[4]
\Needspace{4\baselineskip}
\textbf{解析}

三视图均为“圆内含正方形”，且正方形边长为 $2$。该组合体可还原为棱长为 $2$ 的正方体内接于球：正方体的三个正投影方向恰好给出三个相同的“圆内含正方形”视图。

\Needspace{5\baselineskip}
正方体的体对角线长为
\nopagebreak[4]
\[\sqrt{2^2+2^2+2^2}=2\sqrt3,\]
\Needspace{5\baselineskip}
它等于外接球的直径，故球半径
\nopagebreak[4]
\[R=\frac{2\sqrt3}{2}=\sqrt3.\]
\Needspace{5\baselineskip}
球的表面积为
\nopagebreak[4]
\[S=4\pi R^2=4\pi\times3=12\pi.\]

填 $12\pi$。

\gkanswer{答案：$12\pi$。三视图均为边长 $2$ 的正方形内接于圆，还原为棱长 $2$ 的正方体内接于球，球半径 $R=\sqrt3$，表面积 $4\pi R^2=12\pi$。}

\end{answer}
\begin{answer}
\gksolhead{13}{}

\textbf{题目}：

如图，在 $\triangle ABC$ 中，已知点 $D$ 在 $BC$ 边上，$AD\perp AC$，$\sin\angle BAC=\dfrac{2\sqrt2}{3}$，$AB=3\sqrt2$，$AD=3$，则 $BD$ 的长为\gkblank{}。
\nopagebreak[4]
\begin{center}% ============================================================
% f2013_fujian_science_q13_triangle.tex
% 2013 福建卷（理工类）第 13 题 三角形示意图（TikZ 矢量重绘）
% A 在上，D 在 BC 上，AD 垂直 AC
% 取 A(0,0), D(0,-3), C(3sqrt2,0), B(-sqrt2,-4)
% 整体旋转 -20 度以贴近原页版式
% ============================================================
\begin{tikzpicture}[
  x=1cm, y=1cm, line width=0.6pt, font=\footnotesize, >=Stealth
]

  \begin{scope}[rotate=-20]
    \coordinate (A) at (0,0);
    \coordinate (B) at (-1.4142,-4);
    \coordinate (C) at (4.2426,0);
    \coordinate (D) at (0,-3);

    \draw (A) -- (B) -- (C) -- cycle;
    \draw (A) -- (D);

    % 直角标记（AD 与 AC）
    \draw (0,-0.38) -- (0.38,-0.38) -- (0.38,0);

    \node[above] at (A) {$A$};
    \node[below left] at (B) {$B$};
    \node[below right] at (C) {$C$};
    \node[below left] at (D) {$D$};
  \end{scope}

\end{tikzpicture}
\end{center}
\nopagebreak[4]
\Needspace{4\baselineskip}
\textbf{解析}

\Needspace{5\baselineskip}
设
\nopagebreak[4]
\[\theta=\angle BAD.\]
\Needspace{5\baselineskip}
因为 $AD\perp AC$，且 $D$ 在 $BC$ 边上、$AD$ 位于 $\angle BAC$ 内部，所以
\nopagebreak[4]
\[\angle BAC=\theta+90^\circ.\]
\Needspace{5\baselineskip}
于是
\nopagebreak[4]
\[\cos\theta=\cos\left(\angle BAC-90^\circ\right)=\sin\angle BAC=\frac{2\sqrt2}{3}.\]

\Needspace{5\baselineskip}
在 $\triangle ABD$ 中，$AB=3\sqrt2$，$AD=3$，夹角 $\theta$ 满足 $\cos\theta=\dfrac{2\sqrt2}{3}$。由余弦定理，
\nopagebreak[4]
\[BD^2=AB^2+AD^2-2\cdot AB\cdot AD\cos\theta
=18+9-2\cdot3\sqrt2\cdot3\cdot\frac{2\sqrt2}{3}.\]
\Needspace{5\baselineskip}
计算乘积项：
\nopagebreak[4]
\[2\cdot3\sqrt2\cdot3\cdot\frac{2\sqrt2}{3}=12\cdot2=24.\]
\Needspace{5\baselineskip}
故
\nopagebreak[4]
\[BD^2=27-24=3,\qquad BD=\sqrt3.\]

填 $\sqrt3$。

\gkanswer{答案：$\sqrt3$。由 $AD\perp AC$ 得 $\cos\angle BAD=\sin\angle BAC=\dfrac{2\sqrt2}{3}$，在 $\triangle ABD$ 中用余弦定理得 $BD^2=18+9-24=3$。}

\end{answer}
\begin{answer}
\gksolhead{14}{}

\textbf{题目}：

椭圆 $\Gamma:\dfrac{x^2}{a^2}+\dfrac{y^2}{b^2}=1$（$a>b>0$）的左右焦点分别为 $F_1$，$F_2$，焦距为 $2c$。若直线 $y=\sqrt3(x+c)$ 与椭圆 $\Gamma$ 的一个交点 $M$ 满足 $\angle MF_1F_2=2\angle MF_2F_1$，则该椭圆的离心率等于\gkblank{}。

\Needspace{4\baselineskip}
\textbf{解析}

\Needspace{5\baselineskip}
直线
\nopagebreak[4]
\[y=\sqrt3(x+c)\]
\Needspace{5\baselineskip}
经过左焦点 $F_1(-c,0)$，斜率为 $\sqrt3$，故它与 $x$ 轴正方向的夹角为 $60^\circ$，即
\nopagebreak[4]
\[\angle MF_1F_2=60^\circ.\]
\Needspace{5\baselineskip}
由题设
\nopagebreak[4]
\[\angle MF_1F_2=2\angle MF_2F_1,\]
\Needspace{5\baselineskip}
得
\nopagebreak[4]
\[\angle MF_2F_1=30^\circ.\]
\Needspace{5\baselineskip}
于是
\nopagebreak[4]
\[\angle F_1MF_2=180^\circ-60^\circ-30^\circ=90^\circ,\]
$\triangle MF_1F_2$ 是以 $F_1F_2=2c$ 为斜边的直角三角形。

\Needspace{5\baselineskip}
在这个直角三角形中，
\nopagebreak[4]
\[MF_1=2c\sin30^\circ=c,\qquad MF_2=2c\cos30^\circ=\sqrt3c.\]
\Needspace{5\baselineskip}
由椭圆定义
\nopagebreak[4]
\[MF_1+MF_2=2a,\]
\Needspace{5\baselineskip}
得
\nopagebreak[4]
\[2a=c+\sqrt3c=(1+\sqrt3)c.\]
\Needspace{5\baselineskip}
故离心率
\nopagebreak[4]
\[e=\frac ca=\frac{2}{1+\sqrt3}=\frac{2(\sqrt3-1)}{(\sqrt3+1)(\sqrt3-1)}
=\frac{2(\sqrt3-1)}2=\sqrt3-1.\]

填 $\sqrt3-1$。

\gkanswer{答案：$\sqrt3-1$。直线斜率 $\sqrt3$ 给出 $\angle MF_1F_2=60^\circ$，从而 $\angle MF_2F_1=30^\circ$，$\triangle MF_1F_2$ 为直角三角形，$MF_1=c$，$MF_2=\sqrt3c$，故 $2a=(1+\sqrt3)c$。}

\end{answer}
\begin{answer}
\gksolhead{15}{}

\textbf{题目}：

\Needspace{5\baselineskip}
当 $x\in\mathbf R$，$|x|<1$ 时，有如下表达式：
\nopagebreak[4]
\[1+x+x^2+\cdots+x^n+\cdots=\frac1{1-x}.\]
\Needspace{5\baselineskip}
两边同时积分得：
\nopagebreak[4]
\[\int_0^{\frac12}1\,\mathrm dx+\int_0^{\frac12}x\,\mathrm dx+\int_0^{\frac12}x^2\,\mathrm dx+\cdots+\int_0^{\frac12}x^n\,\mathrm dx+\cdots=\int_0^{\frac12}\frac1{1-x}\,\mathrm dx,\]
\Needspace{5\baselineskip}
从而得到如下等式：
\nopagebreak[4]
\[1\times\frac12+\frac12\times\left(\frac12\right)^2+\frac13\times\left(\frac12\right)^3+\cdots+\frac1{n+1}\times\left(\frac12\right)^{n+1}+\cdots=\ln2.\]
\Needspace{5\baselineskip}
请根据以上材料所蕴含的数学思想方法，计算：
\nopagebreak[4]
\[\mathrm C_n^0\times\frac12+\frac12\mathrm C_n^1\times\left(\frac12\right)^2+\frac13\mathrm C_n^2\times\left(\frac12\right)^3+\cdots+\frac1{n+1}\mathrm C_n^n\times\left(\frac12\right)^{n+1}=\underline{\qquad\qquad}.\]

\Needspace{4\baselineskip}
\textbf{解析}

\Needspace{5\baselineskip}
题中材料给出
\nopagebreak[4]
\[\int_0^{\frac12}x^k\,\mathrm dx=\frac{1}{k+1}\left(\frac12\right)^{k+1}.\]
\Needspace{5\baselineskip}
利用组合恒等式
\nopagebreak[4]
\[\frac{\mathrm C_n^k}{k+1}=\frac{1}{n+1}\mathrm C_{n+1}^{k+1},\]
\Needspace{5\baselineskip}
所求式子可以写成
\nopagebreak[4]
\[\begin{aligned}
&\mathrm C_n^0\cdot\frac12+\frac12\mathrm C_n^1\left(\frac12\right)^2
+\frac13\mathrm C_n^2\left(\frac12\right)^3+\cdots
+\frac1{n+1}\mathrm C_n^n\left(\frac12\right)^{n+1}\\
&=\sum_{k=0}^{n}\frac{\mathrm C_n^k}{k+1}\left(\frac12\right)^{k+1}
=\frac1{n+1}\sum_{k=0}^{n}\mathrm C_{n+1}^{k+1}\left(\frac12\right)^{k+1}.
\end{aligned}\]
\Needspace{5\baselineskip}
令 $j=k+1$，则 $j=1,2,\cdots,n+1$，上式化为
\nopagebreak[4]
\[\frac1{n+1}\sum_{j=1}^{n+1}\mathrm C_{n+1}^{j}\left(\frac12\right)^j.\]
\Needspace{5\baselineskip}
由二项式定理，
\nopagebreak[4]
\[\sum_{j=0}^{n+1}\mathrm C_{n+1}^{j}\left(\frac12\right)^j
=\left(1+\frac12\right)^{n+1}=\left(\frac32\right)^{n+1}.\]
\Needspace{5\baselineskip}
去掉 $j=0$ 的项 $\mathrm C_{n+1}^0=1$，得
\nopagebreak[4]
\[\sum_{j=1}^{n+1}\mathrm C_{n+1}^{j}\left(\frac12\right)^j
=\left(\frac32\right)^{n+1}-1.\]
\Needspace{5\baselineskip}
所以原式
\nopagebreak[4]
\[=\frac1{n+1}\left[\left(\frac32\right)^{n+1}-1\right]
=\frac{3^{n+1}-2^{n+1}}{(n+1)2^{n+1}}.\]

填 $\dfrac{3^{n+1}-2^{n+1}}{(n+1)2^{n+1}}$。

\gkanswer{答案：$\dfrac{3^{n+1}-2^{n+1}}{(n+1)2^{n+1}}$。由 $\dfrac{\mathrm C_n^k}{k+1}=\dfrac{\mathrm C_{n+1}^{k+1}}{n+1}$ 及二项式定理求和可得。}

\end{answer}
\begin{answer}
\gksolhead{16}{}

\textbf{题目}：

某联欢晚会举行抽奖活动，举办方设置了甲、乙两种抽奖方案，方案甲的中奖率为 $\dfrac23$，中奖可以获得 $2$ 分；方案乙的中奖率为 $\dfrac25$，中奖可以获得 $3$ 分；未中奖则不得分。每人有且只有一次抽奖机会，每次抽奖中奖与否互不影响，晚会结束后凭分数兑换奖品。

（1） 若小明选择方案甲抽奖，小红选择方案乙抽奖，记他们的累计得分为 $X$，求 $X\leqslant3$ 的概率；

（2） 若小明、小红两人都选择方案甲或都选择方案乙进行抽奖，问：他们选择何种方案抽奖，累计得分的数学期望较大？

\Needspace{4\baselineskip}
\textbf{解析}

\Needspace{5\baselineskip}
设小明中奖为事件 $A$，小红中奖为事件 $B$，则
\nopagebreak[4]
\[P(A)=\frac23,\qquad P(B)=\frac25,\]
且 $A$，$B$ 相互独立。

\textbf{（1）}

\Needspace{5\baselineskip}
小明得 $2$ 分当且仅当 $A$ 发生，小红得 $3$ 分当且仅当 $B$ 发生，故
\nopagebreak[4]
\[X=2I_A+3I_B.\]
\Needspace{5\baselineskip}
$X\leqslant3$ 的对立事件是 $A$ 与 $B$ 同时发生，此时 $X=5>3$。因此
\nopagebreak[4]
\[P(X\leqslant3)=1-P(A\cap B)=1-P(A)P(B)
=1-\frac23\cdot\frac25=1-\frac4{15}=\frac{11}{15}.\]

\textbf{（2）}

\Needspace{5\baselineskip}
若两人都选择方案甲，则每人的得分期望为
\nopagebreak[4]
\[2\times\frac23=\frac43,\]
\Needspace{5\baselineskip}
两人累计得分的期望为
\nopagebreak[4]
\[E_1=2\times\frac43=\frac83.\]

\Needspace{5\baselineskip}
若两人都选择方案乙，则每人的得分期望为
\nopagebreak[4]
\[3\times\frac25=\frac65,\]
\Needspace{5\baselineskip}
两人累计得分的期望为
\nopagebreak[4]
\[E_2=2\times\frac65=\frac{12}{5}.\]

\Needspace{5\baselineskip}
比较：
\nopagebreak[4]
\[\frac83=\frac{40}{15},\qquad\frac{12}{5}=\frac{36}{15},\]
\Needspace{5\baselineskip}
故
\nopagebreak[4]
\[E_1>E_2.\]
所以两人都选择方案甲抽奖时，累计得分的数学期望较大。

\gkanswer{答案（1）：$\dfrac{11}{15}$。只有两人都中奖时 $X=5>3$，故 $P(X\leqslant3)=1-\dfrac23\times\dfrac25=\dfrac{11}{15}$。}

\gkanswer{答案（2）：都选择方案甲。都选甲时累计得分期望为 $2\times\dfrac23\times2=\dfrac83$；都选乙时为 $2\times\dfrac25\times3=\dfrac{12}{5}$，且 $\dfrac83>\dfrac{12}{5}$。}

\end{answer}
\begin{answer}
\gksolhead{17}{}

\textbf{题目}：

已知函数 $f(x)=x-a\ln x$（$a\in\mathbf R$）。

（1） 当 $a=2$ 时，求曲线 $y=f(x)$ 在点 $A(1,f(1))$ 处的切线方程；

（2） 求函数 $f(x)$ 的极值。

\Needspace{4\baselineskip}
\textbf{解析}

\Needspace{5\baselineskip}
函数定义域为 $x>0$。求导：
\nopagebreak[4]
\[f'(x)=1-\frac ax=\frac{x-a}{x}.\]

\textbf{（1）}

\Needspace{5\baselineskip}
当 $a=2$ 时，
\nopagebreak[4]
\[f(x)=x-2\ln x,\qquad f'(x)=1-\frac2x.\]
\Needspace{5\baselineskip}
在 $x=1$ 处，
\nopagebreak[4]
\[f(1)=1-2\ln1=1,\qquad f'(1)=1-2=-1.\]
\Needspace{5\baselineskip}
切线方程为
\nopagebreak[4]
\[y-1=-1\cdot(x-1),\]
\Needspace{5\baselineskip}
即
\nopagebreak[4]
\[y=-x+2.\]

\textbf{（2）}

\Needspace{5\baselineskip}
由
\nopagebreak[4]
\[f'(x)=\frac{x-a}{x},\qquad x>0,\]
分情况讨论。

若 $a\leqslant0$，则对一切 $x>0$ 都有 $x-a>0$，从而 $f'(x)>0$，$f(x)$ 在 $(0,+\infty)$ 上单调递增，没有极值。

若 $a>0$，则：

\begin{itemize}
\item 当 $0<x<a$ 时，$x-a<0$，$f'(x)<0$，$f(x)$ 单调递减；
\item 当 $x>a$ 时，$x-a>0$，$f'(x)>0$，$f(x)$ 单调递增。
\end{itemize}

\Needspace{5\baselineskip}
所以 $f(x)$ 在 $x=a$ 处取得极小值
\nopagebreak[4]
\[f(a)=a-a\ln a,\]
无极大值。

综上：$a\leqslant0$ 时无极值；$a>0$ 时极小值为 $a-a\ln a$，无极大值。

\gkanswer{答案（1）：$y=-x+2$。$f(1)=1$，$f'(1)=1-2=-1$，切线为 $y-1=-(x-1)$。}

\gkanswer{答案（2）：$a\leqslant0$ 时无极值；$a>0$ 时在 $x=a$ 处取极小值 $a-a\ln a$，无极大值。由 $f'(x)=\dfrac{x-a}{x}$ 讨论即得。}

\end{answer}
\begin{answer}
\gksolhead{18}{}

\textbf{题目}：

如图，在正方形 $OABC$ 中，$O$ 为坐标原点，点 $A$ 的坐标为 $(10,0)$，点 $C$ 的坐标为 $(0,10)$。分别将线段 $OA$ 和 $AB$ 十等分，分点分别记为 $A_1,A_2,\cdots,A_9$ 和 $B_1,B_2,\cdots,B_9$，连接 $OB_i$，过 $A_i$ 作 $x$ 轴的垂线与 $OB_i$ 交于点 $P_i$（$i\in\mathbf N^*$，$1\leqslant i\leqslant9$）。

（1） 求证：点 $P_i$（$i\in\mathbf N^*$，$1\leqslant i\leqslant9$）都在同一条抛物线上，并求该抛物线 $E$ 的方程；

（2） 过点 $C$ 作直线 $l$ 与抛物线 $E$ 交于不同的两点 $M$，$N$，若 $\triangle OCM$ 与 $\triangle OCN$ 的面积之比为 $4:1$，求直线 $l$ 的方程。
\nopagebreak[4]
\begin{center}% ============================================================
% f2013_fujian_science_q18_parabola.tex
% 2013 福建卷（理工类）第 18 题 正方形与抛物线构造（TikZ 矢量重绘）
% O(0,0), A(10,0), B(10,10), C(0,10)
% OA、AB 各十等分；A_i(i,0)，B_i(10,i)；OB_i 与过 A_i 的竖线交于 P_i
% P_i(i, i^2/10) 在抛物线 y=x^2/10 上
% ============================================================
\begin{tikzpicture}[
  x=0.42cm, y=0.42cm, line width=0.6pt, font=\footnotesize, >=Stealth
]

  % ---------- 正方形 ----------
  \draw (0,0) -- (10,0) -- (10,10) -- (0,10) -- cycle;

  % ---------- 十等分刻度 ----------
  \foreach \i in {1,...,9} {
    \draw (\i,0) -- (\i,-0.22);
    \draw (10,\i) -- (10.22,\i);
  }

  % ---------- OB_i 与竖线 ----------
  \foreach \i in {1,...,9} {
    \draw[gray!65, line width=0.35pt] (0,0) -- (10,\i);
    \draw[dashed, gray!65, line width=0.35pt] (\i,0) -- (\i,{\i*\i/10});
    \fill (\i,{\i*\i/10}) circle (0.09);
  }

  % ---------- 标注 ----------
  \node[below left]  at (0,0)   {$O$};
  \node[below right] at (10,0)  {$A$};
  \node[above right] at (10,10) {$B$};
  \node[above left]  at (0,10)  {$C$};

  \node[below] at (1,-0.28) {$A_1$};
  \node[below] at (2,-0.28) {$A_2$};
  \node[below] at (5,-0.28) {$A_i$};
  \node[below] at (9,-0.28) {$A_9$};

  \node[right] at (10.28,1) {$B_1$};
  \node[right] at (10.28,2) {$B_2$};
  \node[right] at (10.28,5) {$B_i$};
  \node[right] at (10.28,9) {$B_9$};

  \node[right] at (5.35,2.95) {$P_i$};

\end{tikzpicture}
\end{center}
\nopagebreak[4]
\Needspace{4\baselineskip}
\textbf{解析}

\textbf{（1）}

\Needspace{5\baselineskip}
由 $OA$ 被十等分，得
\nopagebreak[4]
\[A_i=(i,0)\qquad(i=1,2,\cdots,9).\]
\Needspace{5\baselineskip}
由 $AB$ 被十等分，得
\nopagebreak[4]
\[B_i=(10,i)\qquad(i=1,2,\cdots,9).\]

\Needspace{5\baselineskip}
直线 $OB_i$ 经过 $O(0,0)$ 与 $B_i(10,i)$，其方程为
\nopagebreak[4]
\[y=\frac{i}{10}x.\]
\Needspace{5\baselineskip}
过 $A_i(i,0)$ 作 $x$ 轴的垂线为 $x=i$，与 $OB_i$ 的交点 $P_i$ 满足
\nopagebreak[4]
\[x=i,\qquad y=\frac{i}{10}\cdot i=\frac{i^2}{10}.\]
\Needspace{5\baselineskip}
故
\nopagebreak[4]
\[P_i\left(i,\frac{i^2}{10}\right).\]
\Needspace{5\baselineskip}
其坐标恒满足
\nopagebreak[4]
\[y=\frac{x^2}{10},\]
\Needspace{5\baselineskip}
所以所有 $P_i$ 都在同一条抛物线
\nopagebreak[4]
\[E:\ y=\frac{x^2}{10}\]
上。

\textbf{（2）}

\Needspace{5\baselineskip}
点 $C$ 为 $(0,10)$。设过 $C$ 的直线 $l$ 的斜率为 $k$，则
\nopagebreak[4]
\[l:\ y=kx+10.\]
（若 $l$ 为 $x=0$，它与抛物线只交于原点一个点，不合题意，故斜率存在。）

\Needspace{5\baselineskip}
与抛物线 $y=\dfrac{x^2}{10}$ 联立：
\nopagebreak[4]
\[\frac{x^2}{10}=kx+10
\Longleftrightarrow x^2-10kx-100=0.\]
\Needspace{5\baselineskip}
设两个交点的横坐标分别为 $x_1$，$x_2$，则
\nopagebreak[4]
\[x_1+x_2=10k,\qquad x_1x_2=-100.\]

\Needspace{5\baselineskip}
三角形 $OCM$ 与 $OCN$ 都以 $OC$ 为底，$OC=10$，高分别为 $|x_1|$ 与 $|x_2|$，故
\nopagebreak[4]
\[S_{\triangle OCM}=\frac12\cdot10\cdot|x_1|=5|x_1|,\]
\nopagebreak[4]
\[S_{\triangle OCN}=\frac12\cdot10\cdot|x_2|=5|x_2|.\]
\Needspace{5\baselineskip}
面积之比为 $4:1$，即
\nopagebreak[4]
\[|x_1|:|x_2|=4:1.\]
\Needspace{5\baselineskip}
又 $x_1x_2=-100<0$，说明 $x_1$，$x_2$ 异号。设 $|x_1|=4|x_2|$，则
\nopagebreak[4]
\[|x_1x_2|=4x_2^2=100,\]
\Needspace{5\baselineskip}
故
\nopagebreak[4]
\[|x_2|=5,\qquad |x_1|=20.\]
\Needspace{5\baselineskip}
于是
\nopagebreak[4]
\[x_1+x_2=\pm(20-5)=\pm15.\]
\Needspace{5\baselineskip}
由 $x_1+x_2=10k$ 得
\nopagebreak[4]
\[k=\pm\frac32.\]

\Needspace{5\baselineskip}
所以直线 $l$ 的方程为
\nopagebreak[4]
\[y=\frac32x+10\quad\text{\text{或}}\quad y=-\frac32x+10.\]

\gkanswer{答案（1）：$E:y=\dfrac{x^2}{10}$。$P_i\left(i,\dfrac{i^2}{10}\right)$ 恒满足该方程。}

\gkanswer{答案（2）：$y=\dfrac32x+10$ 或 $y=-\dfrac32x+10$。设 $l:y=kx+10$，由面积比得交点横坐标绝对值之比 $4:1$，结合 $x_1x_2=-100$ 得 $x_1+x_2=\pm15$，故 $k=\pm\dfrac32$。}

\end{answer}
\begin{answer}
\gksolhead{19}{}

\textbf{题目}：

如图，在四棱柱 $ABCD-A_1B_1C_1D_1$ 中，侧棱 $AA_1\perp$ 底面 $ABCD$，$AB\parallel DC$，$AA_1=1$，$AB=3k$，$AD=4k$，$BC=5k$，$DC=6k$（$k>0$）。

（1） 求证：$CD\perp$ 平面 $ADD_1A_1$；

（2） 若直线 $AA_1$ 与平面 $AB_1C$ 所成角的正弦值为 $\dfrac67$，求 $k$ 的值；

（3） 现将与四棱柱 $ABCD-A_1B_1C_1D_1$ 形状和大小完全相同的两个四棱柱拼成一个新的四棱柱，规定：若拼成的新四棱柱形状和大小完全相同，则视为同一种拼接方案。问：共有几种不同的拼接方案？在这些拼接成的新四棱柱中，记其中最小的表面积为 $f(k)$，写出 $f(k)$ 的解析式。（直接写出答案，不必说明理由）
\nopagebreak[4]
\begin{center}% ============================================================
% f2013_fujian_science_q19_prism.tex
% 2013 福建卷（理工类）第 19 题 四棱柱（TikZ 矢量重绘）
% 侧棱 AA1 垂直底面 ABCD；AB 平行 DC
% AB=3k, AD=4k, BC=5k, DC=6k, AA1=1
% 图中另绘平面 AB1C 的三边 AB1、B1C、AC
% ============================================================
\begin{tikzpicture}[
  x=1cm, y=1cm, line width=0.6pt, font=\footnotesize, >=Stealth
]

  % ---------- 顶点（斜二测投影） ----------
  \coordinate (A)  at (0,0);
  \coordinate (B)  at (3,0);
  \coordinate (C)  at (8,1.30);
  \coordinate (D)  at (2,1.30);
  \coordinate (A1) at (0,1.15);
  \coordinate (B1) at (3,1.15);
  \coordinate (C1) at (8,2.45);
  \coordinate (D1) at (2,2.45);

  % ---------- 底面 ----------
  \draw (A) -- (B) -- (C);
  \draw[dashed] (C) -- (D) -- (A);

  % ---------- 侧棱 ----------
  \draw (A) -- (A1);
  \draw (B) -- (B1);
  \draw (C) -- (C1);
  \draw[dashed] (D) -- (D1);

  % ---------- 上底面 ----------
  \draw (A1) -- (B1) -- (C1) -- (D1) -- cycle;

  % ---------- 平面 AB1C ----------
  \draw (A) -- (B1);
  \draw (B1) -- (C);
  \draw[dashed] (A) -- (C);

  % ---------- 标注 ----------
  \node[below left]  at (A)  {$A$};
  \node[below]       at (B)  {$B$};
  \node[right]       at (C)  {$C$};
  \node[above right] at (D)  {$D$};
  \node[left]        at (A1) {$A_1$};
  \node[above]       at (B1) {$B_1$};
  \node[above right] at (C1) {$C_1$};
  \node[above]       at (D1) {$D_1$};

\end{tikzpicture}
\end{center}
\nopagebreak[4]
\Needspace{4\baselineskip}
\textbf{解析}

\textbf{（1）}

\Needspace{5\baselineskip}
先判定底面 $ABCD$ 的形状。由 $AB\parallel DC$，设 $AD$ 与 $DC$ 的夹角为 $\theta$。在底面内取 $D$ 为原点、$DC$ 所在直线为 $x$ 轴，则
\nopagebreak[4]
\[C=(6k,0),\qquad A=(4k\cos\theta,4k\sin\theta).\]
\Needspace{5\baselineskip}
又 $AB\parallel DC$ 且 $AB=3k$，故
\nopagebreak[4]
\[B=A+(3k,0)=(3k+4k\cos\theta,4k\sin\theta).\]
\Needspace{5\baselineskip}
于是
\nopagebreak[4]
\[\begin{aligned}
BC^2&=\left(6k-3k-4k\cos\theta\right)^2+\left(0-4k\sin\theta\right)^2\\
&=k^2\left[(3-4\cos\theta)^2+16\sin^2\theta\right]\\
&=k^2\left(9-24\cos\theta+16\cos^2\theta+16\sin^2\theta\right)\\
&=k^2(25-24\cos\theta).
\end{aligned}\]
\Needspace{5\baselineskip}
已知 $BC=5k$，故
\nopagebreak[4]
\[25k^2=k^2(25-24\cos\theta),\]
\Needspace{5\baselineskip}
即 $\cos\theta=0$。因此
\nopagebreak[4]
\[AD\perp DC.\]
又 $AB\parallel DC$，所以 $AD\perp AB$。

\Needspace{5\baselineskip}
因为 $AA_1\perp$ 底面 $ABCD$，而 $CD\subset$ 底面 $ABCD$，所以
\nopagebreak[4]
\[AA_1\perp CD.\]
\Needspace{5\baselineskip}
在平面 $ADD_1A_1$ 内，$AD$ 与 $AA_1$ 相交于 $A$，且
\nopagebreak[4]
\[CD\perp AD,\qquad CD\perp AA_1.\]
\Needspace{5\baselineskip}
由线面垂直的判定定理，
\nopagebreak[4]
\[CD\perp\text{\text{平面} }ADD_1A_1.\]

\textbf{（2）}

\Needspace{5\baselineskip}
建立空间直角坐标系：以 $A$ 为原点，$AB$ 所在直线为 $x$ 轴，$AD$ 所在直线为 $y$ 轴，$AA_1$ 所在直线为 $z$ 轴。则
\nopagebreak[4]
\[\begin{aligned}
&A(0,0,0),\quad B(3k,0,0),\quad D(0,4k,0),\quad C(6k,4k,0),\\
&A_1(0,0,1),\quad B_1(3k,0,1).
\end{aligned}\]
\Needspace{5\baselineskip}
平面 $AB_1C$ 经过 $A$，$B_1$，$C$，其法向量可取
\nopagebreak[4]
\[\begin{aligned}
\boldsymbol n&=\overrightarrow{AB_1}\times\overrightarrow{AC}\\
&=(3k,0,1)\times(6k,4k,0)\\
&=\left(0\cdot0-1\cdot4k,\ 1\cdot6k-3k\cdot0,\ 3k\cdot4k-0\cdot6k\right)\\
&=(-4k,6k,12k^2).
\end{aligned}\]
\Needspace{5\baselineskip}
直线 $AA_1$ 的方向向量为
\nopagebreak[4]
\[\overrightarrow{AA_1}=(0,0,1).\]
\Needspace{5\baselineskip}
设 $AA_1$ 与平面 $AB_1C$ 所成角为 $\alpha$，则
\nopagebreak[4]
\[\sin\alpha=\frac{\left|\boldsymbol n\cdot\overrightarrow{AA_1}\right|}{\left|\boldsymbol n\right|\left|\overrightarrow{AA_1}\right|}
=\frac{12k^2}{2k\sqrt{13+36k^2}}
=\frac{6k}{\sqrt{13+36k^2}}.\]
\Needspace{5\baselineskip}
已知 $\sin\alpha=\dfrac67$，故
\nopagebreak[4]
\[\frac{6k}{\sqrt{13+36k^2}}=\frac67.\]
\Needspace{5\baselineskip}
两边平方：
\nopagebreak[4]
\[\frac{36k^2}{13+36k^2}=\frac{36}{49}.\]
\Needspace{5\baselineskip}
约去 $36$：
\nopagebreak[4]
\[\frac{k^2}{13+36k^2}=\frac1{49}
\Longleftrightarrow49k^2=13+36k^2
\Longleftrightarrow13k^2=13.\]
\Needspace{5\baselineskip}
因为 $k>0$，所以
\nopagebreak[4]
\[k=1.\]

\textbf{（3）}

\Needspace{5\baselineskip}
先算原四棱柱的底面积与底面周长。由（1）底面为直角梯形，
\nopagebreak[4]
\[S_{\text{\text{底}}}=\frac12(AB+DC)\cdot AD=\frac12(3k+6k)\cdot4k=18k^2,\]
\Needspace{5\baselineskip}
底面周长
\nopagebreak[4]
\[AB+BC+CD+DA=3k+5k+6k+4k=18k.\]
原四棱柱的高为 $AA_1=1$。

两个全等四棱柱拼成新四棱柱，可按拼合面分类。

\Needspace{5\baselineskip}
\textbf{方案一：两底面重合。} 新四棱柱以原梯形为底、高为 $2$，表面积
\nopagebreak[4]
\[S_1=2S_{\text{\text{底}}}+(\text{\text{底面周长}})\times2
=36k^2+36k.\]

\Needspace{5\baselineskip}
\textbf{方案二：沿侧面 $ADD_1A_1$（$4k\times1$）拼接，使 $A$ 与 $D$ 对应。} 把第二个四棱柱关于 $AD$ 所在直线作对称摆放，两个梯形拼成等腰梯形：上底 $12k$、下底 $6k$、高 $4k$、两腰各 $5k$。新底面积
\nopagebreak[4]
\[S'_{\text{\text{底}}}=2\times18k^2=36k^2,\]
\Needspace{5\baselineskip}
周长
\nopagebreak[4]
\[12k+6k+5k+5k=28k.\]
\Needspace{5\baselineskip}
新四棱柱高仍为 $1$，故
\nopagebreak[4]
\[S_2=2\times36k^2+28k=72k^2+28k.\]

\Needspace{5\baselineskip}
\textbf{方案三：沿同一侧面 $ADD_1A_1$ 拼接，但使 $A$ 与 $D$ 互换。} 即把第二个四棱柱绕 $AD$ 的中点旋转 $180^\circ$ 后拼接，两个梯形拼成平行四边形，边长分别为 $9k$ 与 $5k$，高为 $4k$。新底面积仍为
\nopagebreak[4]
\[36k^2,\]
\Needspace{5\baselineskip}
周长
\nopagebreak[4]
\[9k+5k+9k+5k=28k.\]
\Needspace{5\baselineskip}
故
\nopagebreak[4]
\[S_3=2\times36k^2+28k=72k^2+28k.\]
此平行四边形的角不是直角，所得棱柱与方案二的等腰梯形棱柱形状不同。

\Needspace{5\baselineskip}
\textbf{方案四：沿侧面 $BCC_1B_1$（$5k\times1$）拼接，使 $B$ 与 $C$ 互换。} 把第二个四棱柱绕 $BC$ 的中点旋转 $180^\circ$ 后拼接，两个梯形拼成矩形，边长分别为 $9k$ 与 $4k$。新四棱柱为长方体，尺寸 $9k\times4k\times1$，表面积
\nopagebreak[4]
\[S_4=2(9k\cdot4k+9k\cdot1+4k\cdot1)=72k^2+26k.\]

沿 $AB$ 或 $CD$ 拼接时，拼合后的底面不是四边形（为五边形或六边形），不能得到四棱柱，故不计入。

\Needspace{5\baselineskip}
因此共有
\nopagebreak[4]
\[4\text{ \text{种}}\]
不同的拼接方案。

\Needspace{5\baselineskip}
比较表面积：
\nopagebreak[4]
\[S_2=S_3=72k^2+28k>72k^2+26k=S_4,\]
\Needspace{5\baselineskip}
故只需比较 $S_1$ 与 $S_4$。令
\nopagebreak[4]
\[36k^2+36k=72k^2+26k
\Longleftrightarrow36k^2-10k=0
\Longleftrightarrow k(18k-5)=0.\]
\Needspace{5\baselineskip}
因为 $k>0$，得
\nopagebreak[4]
\[k=\frac5{18}.\]
\Needspace{5\baselineskip}
所以
\nopagebreak[4]
\[f(k)=\begin{cases}
36k^2+36k,&0<k\leqslant\dfrac5{18},\\[6pt]
72k^2+26k,&k\geqslant\dfrac5{18}.
\end{cases}\]

\gkanswer{答案（1）：$CD\perp$ 平面 $ADD_1A_1$，见解析正文。由三边长可判 $AD\perp DC$，又 $AA_1\perp$ 底面，故 $CD$ 垂直于平面内两条相交直线 $AD$、$AA_1$。}

\gkanswer{答案（2）：$k=1$。平面 $AB_1C$ 的法向量可取 $(-4k,6k,12k^2)$，由 $\dfrac{6k}{\sqrt{13+36k^2}}=\dfrac67$ 解得 $k=1$。}

\gkanswer{答案（3）：共有 $4$ 种不同的拼接方案；$f(k)=36k^2+36k$（$0<k\leqslant\dfrac5{18}$），$f(k)=72k^2+26k$（$k\geqslant\dfrac5{18}$）。四种方案的表面积分别为 $36k^2+36k$，$72k^2+28k$，$72k^2+28k$，$72k^2+26k$，比较即得。}

\end{answer}
\begin{answer}
\gksolhead{20}{}

\textbf{题目}：

已知函数 $f(x)=\sin(\omega x+\varphi)$（$\omega>0$，$0<\varphi<\pi$）的周期为 $\pi$，图象的一个对称中心为 $\left(\dfrac\pi4,0\right)$。将函数 $f(x)$ 图象上所有点的横坐标伸长到原来的 $2$ 倍（纵坐标不变），再将得到的图象向右平移 $\dfrac\pi2$ 个单位长度后得到函数 $g(x)$ 的图象。

（1） 求函数 $f(x)$ 与 $g(x)$ 的解析式；

（2） 是否存在 $x_0\in\left(\dfrac\pi6,\dfrac\pi4\right)$，使得 $f(x_0)$，$g(x_0)$，$f(x_0)g(x_0)$ 按照某种顺序成等差数列？若存在，请确定 $x_0$ 的个数；若不存在，说明理由；

（3） 求实数 $a$ 与正整数 $n$，使得 $F(x)=f(x)+ag(x)$ 在 $(0,n\pi)$ 内恰有 $2013$ 个零点。

\Needspace{4\baselineskip}
\textbf{解析}

\textbf{（1）}

\Needspace{5\baselineskip}
由周期 $T=\pi$ 及 $\omega>0$，得
\nopagebreak[4]
\[\frac{2\pi}{\omega}=\pi\Longrightarrow\omega=2.\]
\Needspace{5\baselineskip}
图象的一个对称中心为 $\left(\dfrac\pi4,0\right)$，故
\nopagebreak[4]
\[f\left(\frac\pi4\right)=0,\]
\Needspace{5\baselineskip}
即
\nopagebreak[4]
\[\sin\left(2\cdot\frac\pi4+\varphi\right)=\sin\left(\frac\pi2+\varphi\right)=\cos\varphi=0.\]
\Needspace{5\baselineskip}
又 $0<\varphi<\pi$，故
\nopagebreak[4]
\[\varphi=\frac\pi2.\]
\Needspace{5\baselineskip}
于是
\nopagebreak[4]
\[f(x)=\sin\left(2x+\frac\pi2\right)=\cos2x.\]

\Needspace{5\baselineskip}
将 $f(x)=\cos2x$ 图象上所有点的横坐标伸长到原来的 $2$ 倍，纵坐标不变，得到
\nopagebreak[4]
\[y=\cos x.\]
\Needspace{5\baselineskip}
再向右平移 $\dfrac\pi2$ 个单位长度，得到
\nopagebreak[4]
\[g(x)=\cos\left(x-\frac\pi2\right)=\sin x.\]
\Needspace{5\baselineskip}
所以
\nopagebreak[4]
\[f(x)=\cos2x,\qquad g(x)=\sin x.\]

\textbf{（2）}

\Needspace{5\baselineskip}
令
\nopagebreak[4]
\[t=\sin x_0.\]
\Needspace{5\baselineskip}
当
\nopagebreak[4]
\[x_0\in\left(\frac\pi6,\frac\pi4\right)\]
\Needspace{5\baselineskip}
时，
\nopagebreak[4]
\[t\in\left(\frac12,\frac{\sqrt2}{2}\right).\]
\Needspace{5\baselineskip}
此时
\nopagebreak[4]
\[f(x_0)=\cos2x_0=1-2\sin^2x_0=1-2t^2,\]
\nopagebreak[4]
\[g(x_0)=\sin x_0=t,\]
\nopagebreak[4]
\[f(x_0)g(x_0)=t(1-2t^2)=t-2t^3.\]

三个数 $1-2t^2$，$t$，$t-2t^3$ 按某种顺序成等差数列，等价于其中某一个数是另外两个数的算术平均。分三种情况。

\Needspace{5\baselineskip}
\textbf{情形一：$2f(x_0)=g(x_0)+f(x_0)g(x_0)$。}
\nopagebreak[4]
\[2(1-2t^2)=t+(t-2t^3).\]
\Needspace{5\baselineskip}
整理：
\nopagebreak[4]
\[2-4t^2=2t-2t^3
\Longleftrightarrow t^3-2t^2-t+1=0.\]
\Needspace{5\baselineskip}
设
\nopagebreak[4]
\[p(t)=t^3-2t^2-t+1.\]
\Needspace{5\baselineskip}
则
\nopagebreak[4]
\[p\left(\frac12\right)=\frac18-\frac12-\frac12+1=\frac18>0,\]
\nopagebreak[4]
\[p\left(\frac{\sqrt2}{2}\right)=\frac{\sqrt2}{4}-1-\frac{\sqrt2}{2}+1=-\frac{\sqrt2}{4}<0.\]
\Needspace{5\baselineskip}
又
\nopagebreak[4]
\[p'(t)=3t^2-4t-1.\]
\Needspace{5\baselineskip}
在 $\left(\dfrac12,\dfrac{\sqrt2}{2}\right)$ 上，
\nopagebreak[4]
\[3t^2<\frac32,\qquad4t>2,\]
\Needspace{5\baselineskip}
故
\nopagebreak[4]
\[p'(t)<\frac32-2-1=-\frac32<0,\]
$p(t)$ 严格递减。由零点存在定理及单调性，$p(t)=0$ 在该区间内恰有一个根。

\Needspace{5\baselineskip}
\textbf{情形二：$2g(x_0)=f(x_0)+f(x_0)g(x_0)$。}
\nopagebreak[4]
\[2t=(1-2t^2)+(t-2t^3).\]
\Needspace{5\baselineskip}
整理：
\nopagebreak[4]
\[2t^3+2t^2+t-1=0.\]
\Needspace{5\baselineskip}
设
\nopagebreak[4]
\[r(t)=2t^3+2t^2+t-1.\]
\Needspace{5\baselineskip}
则
\nopagebreak[4]
\[r'(t)=6t^2+4t+1>0\qquad(t>0),\]
\Needspace{5\baselineskip}
$r(t)$ 在 $(0,+\infty)$ 上严格递增。又
\nopagebreak[4]
\[r\left(\frac12\right)=\frac14+\frac12+\frac12-1=\frac14>0,\]
故在 $\left(\dfrac12,\dfrac{\sqrt2}{2}\right)$ 上恒有 $r(t)>0$，无解。

\Needspace{5\baselineskip}
\textbf{情形三：$2f(x_0)g(x_0)=f(x_0)+g(x_0)$。}
\nopagebreak[4]
\[2(t-2t^3)=(1-2t^2)+t.\]
\Needspace{5\baselineskip}
整理：
\nopagebreak[4]
\[4t^3-2t^2-t+1=0.\]
\Needspace{5\baselineskip}
设
\nopagebreak[4]
\[q(t)=4t^3-2t^2-t+1.\]
\Needspace{5\baselineskip}
则
\nopagebreak[4]
\[q'(t)=12t^2-4t-1=(2t-1)(6t+1).\]
\Needspace{5\baselineskip}
当 $t>\dfrac12$ 时，$2t-1>0$，$6t+1>0$，故 $q'(t)>0$，$q(t)$ 严格递增。又
\nopagebreak[4]
\[q\left(\frac12\right)=\frac12-\frac12-\frac12+1=\frac12>0,\]
故在 $\left(\dfrac12,\dfrac{\sqrt2}{2}\right)$ 上恒有 $q(t)>0$，无解。

\Needspace{5\baselineskip}
综上，只有情形一给出一个解，即存在 $x_0$，且 $x_0$ 的个数为
\nopagebreak[4]
\[1.\]

\textbf{（3）}

\Needspace{5\baselineskip}
由（1），
\nopagebreak[4]
\[F(x)=f(x)+ag(x)=\cos2x+a\sin x.\]
\Needspace{5\baselineskip}
利用 $\cos2x=1-2\sin^2x$，令
\nopagebreak[4]
\[t=\sin x,\]
\Needspace{5\baselineskip}
则 $F(x)=0$ 等价于
\nopagebreak[4]
\[1-2t^2+at=0,\]
\Needspace{5\baselineskip}
即
\nopagebreak[4]
\[2t^2-at-1=0.\]
\Needspace{5\baselineskip}
其判别式
\nopagebreak[4]
\[\Delta=a^2+8>0,\]
\Needspace{5\baselineskip}
两根为
\nopagebreak[4]
\[t_1=\frac{a+\sqrt{a^2+8}}4>0,\qquad
t_2=\frac{a-\sqrt{a^2+8}}4<0,\]
\Needspace{5\baselineskip}
且
\nopagebreak[4]
\[t_1t_2=-\frac12.\]
所以一个根为正、一个根为负。

记 $N_n$ 为 $F(x)$ 在 $(0,n\pi)$ 内的零点个数。对 $t\in(-1,1)$，方程 $\sin x=t$ 在每个长度为 $2\pi$ 的周期内有两个解；对 $t=\pm1$，每个周期内有一个解。按 $a$ 的取值分类：

\begin{itemize}
\item 若 $a>1$，则 $t_1>1$，$t_2\in(-1,0)$。每个 $2\pi$ 周期内有 $2$ 个零点，故 $N_{2m}=2m$，$N_{2m+1}=2m$，均为偶数。
\item 若 $a=1$，则 $t_1=1$，$t_2=-\dfrac12$。每个 $2\pi$ 周期内，$t=1$ 贡献 $1$ 个零点，$t=-\dfrac12$ 贡献 $2$ 个零点，共 $3$ 个。故
\end{itemize}
\nopagebreak[4]
\[N_{2m}=3m,\qquad N_{2m+1}=3m+1.\]
\begin{itemize}
\item 若 $-1<a<1$，则 $t_1,t_2\in(-1,1)$，每个 $2\pi$ 周期内有 $4$ 个零点，故 $N_{2m}=4m$，$N_{2m+1}=4m+2$，均为偶数。
\item 若 $a=-1$，则 $t_1=\dfrac12$，$t_2=-1$。每个 $2\pi$ 周期内，$t=\dfrac12$ 贡献 $2$ 个零点，$t=-1$ 贡献 $1$ 个零点，共 $3$ 个。故
\end{itemize}
\nopagebreak[4]
\[N_{2m}=3m,\qquad N_{2m+1}=3m+2.\]
\begin{itemize}
\item 若 $a<-1$，则 $t_1\in(0,1)$，$t_2<-1$，每个 $2\pi$ 周期内有 $2$ 个零点，故 $N_{2m}=2m$，$N_{2m+1}=2m+2$，均为偶数。
\end{itemize}

\Needspace{5\baselineskip}
要求 $N_n=2013$。由于 $2013$ 为奇数，只有 $a=1$ 或 $a=-1$ 的情形可能。解
\nopagebreak[4]
\[3m=2013\Longrightarrow m=671,\]
\Needspace{5\baselineskip}
故
\nopagebreak[4]
\[n=2m=1342.\]
\Needspace{5\baselineskip}
而
\nopagebreak[4]
\[3m+1=2013,\qquad3m+2=2013\]
均无整数解（$2012$、$2011$ 都不被 $3$ 整除）。

\Needspace{5\baselineskip}
所以
\nopagebreak[4]
\[a=1\text{ \text{或} }a=-1,\qquad n=1342.\]

\gkanswer{答案（1）：$f(x)=\cos2x$，$g(x)=\sin x$。由周期 $\pi$ 得 $\omega=2$，由对称中心得 $\varphi=\dfrac\pi2$，再作伸缩与平移。}

\gkanswer{答案（2）：存在，且 $x_0$ 的个数为 $1$。令 $t=\sin x_0\in\left(\dfrac12,\dfrac{\sqrt2}{2}\right)$，三种等差顺序中只有 $2f=g+fg$ 对应的方程在区间内有一个根。}

\gkanswer{答案（3）：$a=1$ 或 $a=-1$，$n=1342$。此时每个 $2\pi$ 周期内恰有 $3$ 个零点，$2013=3\times671$，故 $n=2\times671=1342$。}

\end{answer}
\begin{answer}
\gksolhead{21}{三选二}

\textbf{题目}三选二：

以下 A、B、C 三题任选两题作答。

\textbf{【A】} 已知直线 $l:ax+y=1$ 在矩阵 $\boldsymbol A=\begin{pmatrix}1&2\\0&1\end{pmatrix}$ 对应的变换作用下变为直线 $l':x+by=1$。

（1） 求实数 $a$，$b$ 的值；

（2） 若点 $P(x_0,y_0)$ 在直线 $l$ 上，且 $\boldsymbol A\begin{pmatrix}x_0\\y_0\end{pmatrix}=\begin{pmatrix}x_0\\y_0\end{pmatrix}$，求点 $P$ 的坐标。

\textbf{【B】} 在平面直角坐标系中，以坐标原点为极点，$x$ 轴的非负半轴为极轴建立极坐标系。已知点 $A$ 的极坐标为 $\left(\sqrt2,\dfrac\pi4\right)$，直线 $l$ 的极坐标方程为 $\rho\cos\left(\theta-\dfrac\pi4\right)=a$，且点 $A$ 在直线 $l$ 上。

（1） 求 $a$ 的值及直线 $l$ 的直角坐标方程；

（2） 圆 $C$ 的参数方程为 $\begin{cases}x=1+\cos\alpha,\\y=\sin\alpha,\end{cases}$（$\alpha$ 为参数），试判断直线 $l$ 与圆 $C$ 的位置关系。

\textbf{【C】} 设不等式 $|x-2|<a$（$a\in\mathbf N^*$）的解集为 $A$，且 $\dfrac32\in A$，$\dfrac12\notin A$。

（1） 求 $a$ 的值；

（2） 求函数 $f(x)=|x+a|+|x-2|$ 的最小值。

\Needspace{4\baselineskip}
\textbf{解析}

\gksubpart{【A】矩阵与变换}

\textbf{（1）}

\Needspace{5\baselineskip}
矩阵
\nopagebreak[4]
\[\boldsymbol A=\begin{pmatrix}1&2\\0&1\end{pmatrix}\]
\Needspace{5\baselineskip}
对应的变换为
\nopagebreak[4]
\[\begin{pmatrix}x\\y\end{pmatrix}\longmapsto
\begin{pmatrix}1&2\\0&1\end{pmatrix}\begin{pmatrix}x\\y\end{pmatrix}
=\begin{pmatrix}x+2y\\y\end{pmatrix}.\]
\Needspace{5\baselineskip}
设像点为 $(X,Y)$，则
\nopagebreak[4]
\[X=x+2y,\qquad Y=y,\]
\Needspace{5\baselineskip}
反解得
\nopagebreak[4]
\[x=X-2Y,\qquad y=Y.\]
\Needspace{5\baselineskip}
直线 $l:ax+y=1$ 上的点满足
\nopagebreak[4]
\[a(X-2Y)+Y=1,\]
\Needspace{5\baselineskip}
即
\nopagebreak[4]
\[aX+(1-2a)Y=1.\]
\Needspace{5\baselineskip}
这正是像直线 $l'$ 的方程。与
\nopagebreak[4]
\[l':X+bY=1\]
\Needspace{5\baselineskip}
比较系数，得
\nopagebreak[4]
\[a=1,\qquad 1-2a=b.\]
\Needspace{5\baselineskip}
所以
\nopagebreak[4]
\[a=1,\qquad b=-1.\]

\textbf{（2）}

\Needspace{5\baselineskip}
条件
\nopagebreak[4]
\[\boldsymbol A\begin{pmatrix}x_0\\y_0\end{pmatrix}=\begin{pmatrix}x_0\\y_0\end{pmatrix}\]
\Needspace{5\baselineskip}
即
\nopagebreak[4]
\[\begin{pmatrix}x_0+2y_0\\y_0\end{pmatrix}=\begin{pmatrix}x_0\\y_0\end{pmatrix}.\]
\Needspace{5\baselineskip}
比较第一行得
\nopagebreak[4]
\[x_0+2y_0=x_0\Longrightarrow y_0=0.\]
\Needspace{5\baselineskip}
又 $P(x_0,y_0)$ 在直线 $l$ 上，而 $l$ 为
\nopagebreak[4]
\[x+y=1,\]
\Needspace{5\baselineskip}
故
\nopagebreak[4]
\[x_0+y_0=1.\]
\Needspace{5\baselineskip}
代入 $y_0=0$，得
\nopagebreak[4]
\[x_0=1.\]
\Needspace{5\baselineskip}
所以
\nopagebreak[4]
\[P(1,0).\]

\gksubpart{【B】坐标系与参数方程}

\textbf{（1）}

\Needspace{5\baselineskip}
点 $A$ 的极坐标为 $\left(\sqrt2,\dfrac\pi4\right)$，化为直角坐标：
\nopagebreak[4]
\[x_A=\sqrt2\cos\frac\pi4=1,\qquad y_A=\sqrt2\sin\frac\pi4=1,\]
即 $A(1,1)$。

\Needspace{5\baselineskip}
直线 $l$ 的极坐标方程为
\nopagebreak[4]
\[\rho\cos\left(\theta-\frac\pi4\right)=a.\]
\Needspace{5\baselineskip}
点 $A$ 在 $l$ 上，代入 $\rho=\sqrt2$，$\theta=\dfrac\pi4$：
\nopagebreak[4]
\[\sqrt2\cos\left(\frac\pi4-\frac\pi4\right)=\sqrt2\cos0=\sqrt2=a.\]
\Needspace{5\baselineskip}
故
\nopagebreak[4]
\[a=\sqrt2.\]

\Needspace{5\baselineskip}
将
\nopagebreak[4]
\[\rho\cos\left(\theta-\frac\pi4\right)=\sqrt2\]
\Needspace{5\baselineskip}
展开：
\nopagebreak[4]
\[\rho\left(\cos\theta\cos\frac\pi4+\sin\theta\sin\frac\pi4\right)=\sqrt2,\]
\Needspace{5\baselineskip}
即
\nopagebreak[4]
\[\frac{x+y}{\sqrt2}=\sqrt2.\]
\Needspace{5\baselineskip}
所以直线 $l$ 的直角坐标方程为
\nopagebreak[4]
\[x+y=2.\]

\textbf{（2）}

\Needspace{5\baselineskip}
圆 $C$ 的参数方程为
\nopagebreak[4]
\[\begin{cases}x=1+\cos\alpha,\\y=\sin\alpha,\end{cases}\]
\Needspace{5\baselineskip}
消去 $\alpha$ 得
\nopagebreak[4]
\[(x-1)^2+y^2=1.\]
圆心为 $(1,0)$，半径为 $1$。

\Needspace{5\baselineskip}
圆心 $(1,0)$ 到直线 $x+y-2=0$ 的距离为
\nopagebreak[4]
\[d=\frac{|1+0-2|}{\sqrt{1^2+1^2}}=\frac{1}{\sqrt2}=\frac{\sqrt2}{2}.\]
\Needspace{5\baselineskip}
因为
\nopagebreak[4]
\[d=\frac{\sqrt2}{2}<1,\]
所以直线 $l$ 与圆 $C$ 相交于两个不同的点。

\gksubpart{【C】不等式选讲}

\textbf{（1）}

\Needspace{5\baselineskip}
不等式
\nopagebreak[4]
\[|x-2|<a\qquad(a\in\mathbf N^*)\]
\Needspace{5\baselineskip}
的解集为
\nopagebreak[4]
\[A=(2-a,2+a).\]

\Needspace{5\baselineskip}
由 $\dfrac32\in A$，得
\nopagebreak[4]
\[\left|\frac32-2\right|=\frac12<a,\]
\Needspace{5\baselineskip}
即
\nopagebreak[4]
\[a>\frac12.\]

\Needspace{5\baselineskip}
由 $\dfrac12\notin A$，得
\nopagebreak[4]
\[\left|\frac12-2\right|=\frac32\geqslant a,\]
\Needspace{5\baselineskip}
即
\nopagebreak[4]
\[a\leqslant\frac32.\]

\Needspace{5\baselineskip}
结合 $a\in\mathbf N^*$，得
\nopagebreak[4]
\[\frac12<a\leqslant\frac32,\qquad a\in\{1,2,3,\cdots\}.\]
\Needspace{5\baselineskip}
故
\nopagebreak[4]
\[a=1.\]

\textbf{（2）}

\Needspace{5\baselineskip}
当 $a=1$ 时，
\nopagebreak[4]
\[f(x)=|x+1|+|x-2|.\]
\Needspace{5\baselineskip}
由绝对值三角不等式，
\nopagebreak[4]
\[f(x)=|x+1|+|x-2|\geqslant|(x+1)-(x-2)|=|3|=3.\]
当 $-1\leqslant x\leqslant2$ 时，$x+1\geqslant0$，$x-2\leqslant0$，故 $f(x)=(x+1)+(2-x)=3$。
\Needspace{5\baselineskip}
所以等号可以取到，最小值为
\nopagebreak[4]
\[3.\]

\gkanswer{答案（1）：A：$a=1$，$b=-1$；B：$a=\sqrt2$，直线 $l:x+y=2$；C：$a=1$。分别由矩阵变换、极坐标代入和绝对值不等式解集条件得到。}

\gkanswer{答案（2）：A：$P(1,0)$；B：直线 $l$ 与圆 $C$ 相交于两个不同点；C：最小值为 $3$。A 由不动点条件与 $P\in l$ 解得；B 圆心到直线距离 $\dfrac{\sqrt2}{2}<1$；C 由三角不等式取等。}

\end{answer}